% TOP_PROBLEM: {"problemId": "problem.brauers-kb-conjecture", "canonicalTitle": "Brauer's k(B) Conjecture", "releaseRank": 352, "primaryDomain": "algebra_representation_category", "primaryDomainLabel": "Algebra, representation & category theory", "displayStatus": "open_with_solved_subcases", "releaseStatus": "open_with_solved_subcases"}
% !TeX program = lualatex
% ATTEMPT_STATUS: unresolved
\documentclass[11pt]{article}
\usepackage[margin=25mm]{geometry}
\usepackage{fontspec}
\setmainfont{Latin Modern Roman}
\usepackage{amsmath,amssymb,mathtools}
\usepackage{unicode-math}
\setmathfont{Latin Modern Math}
\usepackage{xurl,hyperref}
\hypersetup{colorlinks=true,urlcolor=blue}
\setcounter{secnumdepth}{0}
\setlength{\parindent}{0pt}
\setlength{\parskip}{5pt}
\setlength{\emergencystretch}{3em}
\begin{document}
\section*{\texorpdfstring{Brauer's k(B) Conjecture}{Brauer's k(B) Conjecture}}\label{352}
\subsection{Definitions and mathematical statement}
Fix a finite group $G$ and a prime $p$. Its ordinary irreducible complex characters split into $p$-blocks via modular representation theory over a splitting $p$-modular system. For a block $B$, write $\operatorname{Irr}(B)$ for its ordinary characters and $D$ for a defect group, defined up to conjugacy by the corresponding maximal Brauer pair. Brauer's conjecture asserts
\[
 k(B):=|\operatorname{Irr}(B)|\le |D|.
\]
The inequality applies to every finite group, every prime, and every block; $D$ need not be abelian. It counts ordinary characters, not the generally different number $\ell(B)$ of irreducible modular representations. No equality criterion or specified correspondence between characters and elements of $D$ is required.
\par\smallskip\textit{Formulation and concept references:} \href{https://arxiv.org/pdf/1911.10710}{[S1]}.

\subsection{Short English statement}
Can the number of ordinary irreducible characters in any block ever exceed the order of its defect group?
\subsection{Research attempt}
\paragraph{Scope and source audit (27 September 2026).}
The quantity is the number of ordinary irreducible characters in one
$p$-block, not the number of all characters of the group and not the
number $\ell(B)$ of simple modular modules. The defect group may be
nonabelian. A proof for blocks of $p$-groups, or even for every abelian
defect group, would not by itself cover the stated scope.

Ardito--Sambale [S1] prove further Cartan-matrix and defect-group cases;
their primary paper was inspected. Sambale's author-hosted book [E1]
provides the basic block-theoretic facts used explicitly below. We do
not treat its list of known cases as a proof of the unrestricted
conjecture. In particular numerical constraints on an arbitrary positive
integral matrix are weaker than being the Cartan matrix of a block.

The approach reconstructs the $p$-group and abelian-group cases,
extends them to direct products with $p'$-groups, proves a general
quadratic-form certificate, and supplies numerical counterexamples to
several overoptimistic matrix reductions. All ordinary character counts
use the established theorem that irreducible complex characters form
a basis of class functions, so their number equals the number of
conjugacy classes.

\paragraph{Why a finite $p$-group has exactly one block.}
Let $P$ be a finite $p$-group and $k$ an algebraically closed field of
characteristic $p$. Write $I$ for the augmentation ideal in $kP$,
the kernel of $\sum a_g g\mapsto\sum a_g$. We prove that $I$ is
nilpotent by induction on $|P|$.

If $P=1$, the ideal is zero. Otherwise the class equation shows that
$Z(P)$ has order divisible by $p$, and hence contains an element $z$
of order $p$. The element $z-1$ is central in the group algebra and
\[
 (z-1)^p=z^p-1=0.
\]
Quotienting by its generated ideal identifies the group algebra with
$k[P/\langle z\rangle]$. The image of $I$ is the augmentation ideal
there, so by induction some $I^a$ is contained in $(z-1)kP$.
Consequently
\[
 I^{ap}\subseteq((z-1)kP)^p=(z-1)^p kP=0.
\]
This proves nilpotence, without claiming that every nil ideal in an
arbitrary infinite-dimensional ring is nilpotent.

An element of $kP$ with nonzero augmentation is a nonzero scalar times
$1+x$ for $x\in I$, and is invertible by a finite geometric series.
Elements of $I$ are nilpotent and cannot be invertible. Hence $kP$ is
local, with unique maximal ideal $I$ and residue field $k$.
An idempotent $e$ has augmentation zero or one. In the first case
$e\in I$ is nilpotent, so $e=0$; in the second $1-e$ is a nilpotent
idempotent, so $e=1$. Thus there is no nontrivial central idempotent.
The decomposition into blocks, defined by primitive central idempotents,
has just one block.

That block is the principal block and has defect group $P$, by the
standard principal-block defect theorem. This is the only defect-group
input needed for the next conclusion.

\paragraph{A complete proof for the blocks of $p$-groups.}
All ordinary irreducible characters of $P$ lie in its unique block $B$.
The character-class correspondence gives
\[
 k(B)=k(P)=\#\{\text{conjugacy classes of }P\}\le |P|=|D|.
 \tag{352.1}
\]
The inequality is just a partition count: every class contains at least
one element. Equality holds in this subcase exactly when every class
is a singleton, equivalently when $P$ is abelian.

For a concrete strict example, the dihedral group of order eight has
five conjugacy classes:
\[
 \{1\},\quad\{r^2\},\quad\{r,r^3\},\quad
 \{s,r^2s\},\quad\{rs,r^3s\}.
\]
Its unique $2$-block has $k(B)=5<8$. The quaternion group of order
eight also has five classes, namely its two central singleton classes
and the three pairs $\{\pm i\},\{\pm j\},\{\pm k\}$.
These examples demonstrate that nonabelian defect groups are already
included in the elementary positive class, although not arbitrary
blocks with those defect groups.

The proof cannot be extended by replacing $P$ with an arbitrary $G$:
the class count only gives $k(B)\le k(G)\le|G|$. The desired quantity
on the right is $|D|$, which may be much smaller than $|G|$.

\paragraph{Every block of a finite abelian group has equality.}
Let $G=P\times H$ be finite abelian, with $P$ its Sylow $p$-subgroup
and $H$ of order prime to $p$. The splitting group algebra $kH$ has
primitive idempotents
\[
 e_\chi=\frac1{|H|}\sum_{h\in H}\chi(h)^{-1}h,
 \qquad \chi\in\operatorname{Hom}(H,k^*).
 \tag{352.2}
\]
Character orthogonality, or direct summation of the roots of unity,
shows $e_\chi e_\psi=0$ for $\chi\ne\psi$, $e_\chi^2=e_\chi$,
and $\sum_\chi e_\chi=1$. Each summand of $kG$ cut out by one
$e_\chi$ is isomorphic to $kP$, and hence is a single block by the
preceding local-algebra argument. There are $|H|$ blocks.

In an ordinary splitting field, every irreducible character is a tensor
product $\theta\otimes\chi$, where $\theta$ is a linear character
of $P$ and $\chi$ a linear character of $H$. Reduction of a root of
unity of $p$-power order is one in characteristic $p$: the polynomial
$X^{p^a}-1$ becomes $(X-1)^{p^a}$. Roots of unity of order prime to
$p$ retain their separate characters in the splitting modular system.
Thus precisely the $|P|$ characters with a fixed $\chi$ lie in the
block labeled by $e_\chi$.

Its defect group is $P$. This can be checked from the Brauer-map
definition: $P$ centralizes all of $G$, so the Brauer map on the
central idempotent (352.2) leaves it nonzero. No larger $p$-subgroup
exists. Therefore every block satisfies
\[
 k(B)=|P|=|D|.
 \tag{352.3}
\]
The argument distinguishes ordinary characters, modular simple modules,
and defect groups rather than treating all three as the same count.
Here $\ell(B)=1$, even though $k(B)$ can be arbitrarily large.

\paragraph{A direct-product extension including nonabelian factors.}
Now let $G=P\times H$, where $P$ is any finite $p$-group and $H$
is any finite $p'$-group. Maschke's theorem and the splitting hypothesis
give
\[
 kH\cong\prod_{\chi\in\operatorname{Irr}(H)}M_{\chi(1)}(k).
\]
We use here the standard characteristic-prime-to-order correspondence:
the ordinary irreducibles of $H$ match the simple modules over the
splitting residue field, with the same degrees. It is a semisimple
case of modular representation theory, not a claim for $p\mid|H|$.
Tensoring with $kP$ decomposes $kG$ into matrix algebras over $kP$.
Their centers have no further nontrivial idempotents, so each is a block.

Ordinary irreducibles of a direct product are the tensor products
$\theta\otimes\chi$. The block corresponding to $\chi$ contains
exactly those with $\theta\in\operatorname{Irr}(P)$, hence has
$k(B)=k(P)$. Its defect group is the Sylow subgroup $P$: its block
idempotent is supported on the commuting $H$ factor and has nonzero
Brauer image at $P$. Thus (352.1) proves the conjecture for every
block in this direct-product class. No claim is made that a group with
normal Sylow subgroup must split as this commuting direct product.
Nontrivial conjugation and fusion can change the block structure.

\paragraph{A general decomposition-matrix inequality.}
For an arbitrary block, let $Q=(d_{\chi\varphi})$ be its ordinary
decomposition matrix, with $k=k(B)$ rows and $l=\ell(B)$ columns.
The standard decomposition theory gives nonnegative integer entries,
no zero row, full column rank, and the Cartan identity
\[
 C=Q^{\mathsf T}Q.
 \tag{352.4}
\]
Every row is nonzero because the restriction of an ordinary character
to the $p$-regular elements has a nonzero value at the identity and
a nonzero expansion in irreducible Brauer characters.

Taking the trace of (352.4) gives the completely elementary bound
\[
 k\le\sum_{\chi,\varphi}d_{\chi\varphi}^2
       =\operatorname{tr}(C).
 \tag{352.5}
\]
This is useful only if the trace itself can be controlled. The fact
that $|D|$ is the largest elementary divisor of $C$ does not mean
that its trace is at most $|D|$, or even that its determinant equals
$|D|$. Elementary divisors are invariant factors under integral row
and column changes, not eigenvalues or diagonal entries.

\paragraph{A complete block-theoretic subcase: one modular simple module.}
The standard Cartan elementary-divisor theorem says that the largest
elementary divisor is $|D|$ and occurs once. If $l=1$, the Cartan
matrix is a one-by-one matrix whose only entry must therefore be $|D|$.
Equation (352.4) becomes
\[
 \sum_{\chi\in\operatorname{Irr}(B)}d_{\chi\varphi}^2=|D|.
\]
Each term is a positive integer, so $k(B)\le|D|$. This proves the
conjecture for all blocks with $\ell(B)=1$, using that explicit
standard block-theoretic input. Equality in this argument forces every
decomposition number in the column to be one. We do not claim an
equality characterization for arbitrary blocks.

This case contains the abelian and direct-product examples above,
but its proof does not require $G$ itself to have a direct-product
structure. It also illustrates how one extra genuine block constraint
can turn the coarse trace bound into exactly the desired inequality.

\paragraph{A quadratic-form certificate with a complete proof.}
\textbf{Lemma 352.1.} Let $A$ be a real symmetric positive-definite
$l\times l$ matrix such that
\[
 v^{\mathsf T}Av\ge1\quad\text{for every }0\ne v\in\mathbb Z^l.
 \tag{352.6}
\]
Then a block with Cartan matrix $C$ satisfies
$k(B)\le\operatorname{tr}(AC)$. In particular,
$\operatorname{tr}(AC)\le|D|$ proves the conjecture for that block.

\emph{Proof.} Apply (352.6) to each nonzero integral row $q_\chi$
of $Q$. Summing gives
\[
 k(B)\le\sum_\chi q_\chi A q_\chi^{\mathsf T}
       =\operatorname{tr}(A Q^{\mathsf T}Q)
       =\operatorname{tr}(AC).
\]
The middle equality follows by expanding the finite sums over matrix
entries, or by cyclicity of the trace.\hfill$\square$

One may first change the basic-set coordinates by an integral unimodular
matrix $U$. Then $Q'=QU$ still has nonzero integral rows and
$C'=U^{\mathsf T}CU$; the same proof applies. Nonnegativity of the
new rows is unnecessary. But the lower bound in (352.6) must hold on
\emph{all} nonzero integer vectors, not only on a finite search box.
A rigorous shortest-vector bound or a direct algebraic expression can
supply that certificate.

For example the form
\[
 q(v_1,\ldots,v_l)=\sum_{i=1}^l v_i^2
                         -\sum_{i=1}^{l-1}v_i v_{i+1}
\]
is positive definite and has minimum at least one on nonzero integer
vectors. Indeed
\[
 2q=v_1^2+v_l^2+\sum_{i=1}^{l-1}(v_i-v_{i+1})^2.
\]
It is positive unless all entries vanish, and $q$ is an integer.
Its matrix has diagonal one and adjacent off-diagonal entries $-1/2$,
so the certificate yields
\[
 k(B)\le\sum_i c_{ii}-\sum_{i=1}^{l-1}c_{i,i+1}.
 \tag{352.7}
\]
This is a genuine improvement over the trace when the adjacent entries
are positive. It does not say that some ordering always makes the
right side at most $|D|$.

\paragraph{Artificial Cartan data expose a missing block constraint.}
For an integer $d\ge2$, take $Q_d$ with $d$ identical rows $(1,1)$
and one row $(0,1)$. It has $d+1$ nonzero nonnegative integral rows,
full column rank, and
\[
 C_d=Q_d^{\mathsf T}Q_d
     =\begin{pmatrix}d&d\\d&d+1\end{pmatrix},
 \qquad\det C_d=d.
 \tag{352.8}
\]
The greatest common divisor of its entries is one, so its Smith
invariant factors are $1,d$. Its off-diagonal entry is positive,
so it is not a direct sum after merely permuting the two coordinates.
Nevertheless its number of rows is $d+1>d$.

These matrices are not claimed to come from blocks of defect order
$d$. They show that (352.4), nonnegative nonzero rows, positive
definiteness, elementary divisors $1,d$, and this elementary notion
of connectedness alone do not prove the target. In fact an integral
change of columns sends the repeated row to $(1,0)$ and the last row
to $(0,1)$, exhibiting stronger decomposability that genuine block
arguments must rule out. Ignoring such extra structure is exactly
why a purely unrestricted matrix proof would be false.

Applied to (352.8), the form in (352.7) gives
$2d+1-d=d+1$, detecting rather than concealing the obstruction.
Trying finitely many quadratic forms on genuine examples is a search
for certificates; a proof that a successful form always exists would
be a new required theorem, not part of Lemma 352.1.

\paragraph{A known determinant bound is not automatically a defect bound.}
For context, Sambale's primary paper [E2] proves a stronger general
bound in terms of $\det C$. We do not reproduce its structural proof.
Even using a determinant bound, one must distinguish
$\det C$ from the largest elementary divisor $|D|$: the determinant
is the product of all elementary divisors. When other factors exceed
one, it may be much larger than $|D|$.

Thus a successful route through Cartan determinants would additionally
need control of those factors, a suitable subsection reduction, or a
sharper form-based estimate exploiting the actual block. Replacing
their product by their maximum is not an inequality in the helpful
direction. The artificial family above is also not a counterexample
to the established determinant theorem for actual blocks, since it
has not met the latter's block hypotheses.

\paragraph{Checks and the exact unresolved step.}
The companion exact script computes conjugacy classes of the displayed
order-eight groups and verifies the group-algebra central nilpotence
identity in small $p$-group cases. It also checks (352.8), its Smith
factors through the entry gcd and determinant, and the quadratic-form
trace identity on explicit integral matrices. The universal lower bound
for the path form is proved by the displayed sum of squares, not
inferred from the finite enumeration.

\textbf{Outcome: UNRESOLVED.} The $p$-group, abelian-group, direct-product
and one-modular-simple cases are proved using the stated basic inputs.
The quadratic-form lemma is a complete sufficient certificate for
individual blocks. The remaining general step is to exploit enough
of the restrictions on genuine decomposition matrices, defects and
fusion to obtain $|D|$ rather than only a trace or determinant bound.
No general certificate construction or classification accomplishing
that step is established here.

\subsection{Sources}
\begin{itemize}
\item[S1]Cesare Giulio Ardito and Benjamin Sambale, "Cartan matrices and Brauer's k(B)-Conjecture V," arXiv:1911.10710 (2019).
\url{https://arxiv.org/pdf/1911.10710}.
\textit{Formulation source.}
\item[E1]Benjamin Sambale, \emph{Blocks of Finite Groups and Their
Invariants}, author-hosted book text.
\url{https://benjaminsambale.github.io/pdfs/lnm.pdf}.
\textit{Inspected 27 September 2026, especially the basic block facts
and Cartan methods. Standard decomposition and defect inputs are
identified explicitly in the attempt.}
\item[E2]Benjamin Sambale, \emph{Cartan matrices and Brauer's
$k(B)$-Conjecture III}, arXiv:1412.7017.
\url{https://arxiv.org/abs/1412.7017}.
\textit{Primary abstract inspected 27 September 2026. Its determinant
bound concerns genuine blocks, not arbitrary nonnegative Gram matrices.}
\end{itemize}
\end{document}
